% ch8_a §8.1–8.2 (书页 255–265 / p269–p279)
\chapter{光学性质}

\epigraphbook{但是，轻些！那是什么光从那边的窗户透出来……}{Romeo and Juliet}

\section{宏观理论}

在本章中，我们讨论电磁波射入固体并在其中传播的问题。众所周知，有些固体是\emph{透明的}，另一些是\emph{不透明的}；有些固体表面\emph{强烈反射}，另一些则倾向于\emph{吸收}照射到其上的辐射。这些效应依赖于辐射的频率；因此我们把整个电磁波谱——从长波无线电波一直到软 X 射线——都包括进来。

电磁波是 Maxwell 方程组的解。我们把它写成如下形式
\begin{equation*}
\left.\begin{aligned}
\nabla\wedge\mathbf{H} &= \frac{\epsilon}{c}\,\frac{\partial\mathbf{E}}{\partial t} + \frac{4\pi\sigma}{c}\,\mathbf{E}, &\qquad \nabla\cdot\mathbf{E} &= 0,\\[4pt]
\nabla\wedge\mathbf{E} &= -\frac{\mu}{c}\,\frac{\partial\mathbf{H}}{\partial t}, &\qquad \nabla\cdot\mathbf{H} &= 0.
\end{aligned}\right\}\tag{8.1}
\end{equation*}
在本章中我们将不考虑任何磁效应：取 $\mu = 1$。

自由空间中的传播同固体中的传播这两者的差别，由两个系数表述——介电常数（dielectric constant）$\epsilon$ 和电导率 $\sigma$。前者给出因 $\mathbf{E}$ 随时间变化而引起的位移电流的大小；后者则是电场所产生的真实电流的一种量度。

从这些方程中消去磁场，得
\begin{equation*}
\nabla^2\mathbf{E} = \frac{\epsilon}{c^2}\,\frac{\partial^2\mathbf{E}}{\partial t^2} + \frac{4\pi\sigma}{c^2}\,\frac{\partial\mathbf{E}}{\partial t}\,.\tag{8.2}
\end{equation*}
这代表一个伴随耗散而传播的波。若选定频率 $\omega$，并写成
\begin{equation*}
\mathbf{E} = \mathbf{E}_0\exp\{i(\mathbf{K}\cdot\mathbf{r}-\omega t)\},\tag{8.3}
\end{equation*}
则波动方程要求
\begin{equation*}
-\mathbf{K}^2 = -\epsilon\,\frac{\omega^2}{c^2} - \frac{4\pi\sigma i\omega}{c^2}\,,\tag{8.4}
\end{equation*}
即
\begin{equation*}
K = \frac{\omega}{c}\left(\epsilon + \frac{4\pi\sigma i}{\omega}\right)^{\frac{1}{2}}.\tag{8.5}
\end{equation*}
这样，一般说来传播常数 $K$ 就是一个复数。而在自由空间中，我们应当简单地有
\begin{equation*}
K = \frac{\omega}{c},\tag{8.6}
\end{equation*}
正如一个以光速行进的波那样。在介质中，速度有所改变：我们说相速被一个\emph{复折射率}（complex refractive index）所除
\begin{equation*}
\mathsf{N} = \left(\epsilon + \frac{4\pi\sigma i}{\omega}\right)^{\frac{1}{2}}.\tag{8.7}
\end{equation*}
从宏观上观察到的光学性质的整个理论，都可以用 $\mathsf{N}$ 来表述。例如，设有一个频率为 $\omega$ 的扰动试图作为平面波沿 $z$ 方向传播，并设我们写出
\begin{equation*}
\mathsf{N} = \mathsf{n} + i\mathsf{k}\tag{8.8}
\end{equation*}
表示 $\mathsf{N}$ 的实部和虚部。于是传播常数变为
\begin{equation*}
K = \frac{\mathsf{n}\omega}{c} + \frac{i\mathsf{k}\omega}{c},\tag{8.9}
\end{equation*}
从而波 (8.3) 变为
\begin{equation*}
\mathbf{E} = \mathbf{E}_0\exp\left\{i\omega\left(\frac{\mathsf{n}z}{c} - t\right)\right\}\exp\left(-\frac{\mathsf{k}\omega z}{c}\right),\tag{8.10}
\end{equation*}
速度减小为 $c/\mathsf{n}$，而波在前进时受到\emph{阻尼}，每行进一个波长就衰减一个比例 $\exp(-2\pi\mathsf{k}/\mathsf{n})$。

波的这种阻尼当然与电磁能量的吸收相联系。为计算吸收，我们应当用 Maxwell 方程求出与 (8.10) 相伴随的电流。它就是 (8.1) 第一式的右端
\begin{equation*}
\begin{aligned}
\mathbf{J} &= \left(-\frac{i\omega\epsilon}{c} + \frac{4\pi\sigma}{c}\right)\mathbf{E}\\
&= -\frac{i\omega}{c}\,\mathsf{N}^2\mathbf{E}
\end{aligned}\tag{8.11}
\end{equation*}
依式 (8.7)。产生焦耳热的速率是下式的实部：
\begin{equation*}
\mathbf{J}\cdot\mathbf{E} = -\frac{i\omega}{c}\,\mathsf{N}^2 E^2.\tag{8.12}
\end{equation*}
于是，\emph{吸收系数}（absorption coefficient）——即每通过单位厚度材料所吸收的能量份额——由下式给出
\begin{equation*}
\eta = \frac{\mathrm{Re}(\mathbf{J}\cdot\mathbf{E})}{\mathsf{n}|\mathbf{E}|^2} = \frac{2\mathsf{k}\omega}{c}\,.\tag{8.13}
\end{equation*}
常被研究的另一种情形，是辐射入射到材料的平面表面上。一般说来，这是一个复杂的光学问题，但我们只需考虑
%FIG{136}: {}
垂直入射的情形。我们要构造 Maxwell 方程组的一个解，使其在介质内部具有 (8.10) 的形式，而在外部则与一个入射波和一个反射波相匹配。于是，对 $z > 0$ 我们写出
\begin{equation*}
E_x = E_0\exp\{i\omega(\mathsf{N}z/c - t)\},\tag{8.14}
\end{equation*}
而对于 $z < 0$，即在外部自由空间中，我们写出
\begin{equation*}
E_x = E_1\exp\{i\omega(z/c - t)\} + E_2\exp\{-i\omega(z/c + t)\},\tag{8.15}
\end{equation*}
分别对应于振幅为 $E_1$、向右行进的波和振幅为 $E_2$、向左行进的波。在边界上使二者匹配，得
\begin{equation*}
E_0 = E_1 + E_2.\tag{8.16}
\end{equation*}
但这些波还伴随有磁场，其磁矢量譬如说沿 $y$ 方向。用 Maxwell 方程组算出 $H_y$，得
\begin{equation*}
-\mathsf{N}E_0 = E_2 - E_1\tag{8.17}
\end{equation*}
此式是与外部的波相匹配而得到的。于是，反射波与入射波的复振幅之比为
\begin{equation*}
\frac{E_2}{E_1} = \frac{1-\mathsf{N}}{1+\mathsf{N}},\tag{8.18}
\end{equation*}
它对应于一个实的\emph{反射系数}（reflection coefficient）
\begin{equation*}
\mathsf{R} = \left|\frac{1-\mathsf{N}}{1+\mathsf{N}}\right|^2 = \frac{(\mathsf{n}-1)^2+\mathsf{k}^2}{(\mathsf{n}+1)^2+\mathsf{k}^2}\tag{8.19}
\end{equation*}
以复折射率 $\mathsf{N}$ 的实部和虚部即式 (8.8) 表出。

显然，对反射系数和吸收系数各作一次独立测量，便足以定出 $\mathsf{n}$ 和 $\mathsf{k}$ 的值。因此，它们就是人们通常引用的光学常数，其理论正是我们将在本章中讨论的。实际中的实验可能要复杂得多——例如相对表面成一定角度的反射，或穿过薄膜的透射——但从原则上讲，其结果都应当由这同样两个系数来描述。

然而，系数 $\mathsf{n}$ 和 $\mathsf{k}$ 并非彼此完全独立。它们由\emph{色散关系}（dispersion relations）联系起来。(8.11) 中的 $\mathsf{N}^2$ 这类量，就是\emph{广义极化率}（generalized susceptibility）$\alpha(\omega)$ 的一个例子，出现在如下关系中
\begin{equation*}
D(\omega) = \alpha(\omega)\,F(\omega)\tag{8.20}
\end{equation*}
它联系着广义的“位移” $D$ 与“力” $F$ 在某个频率 $\omega$ 处的 Fourier 分量。正如初等交流电路理论中那样，它必须是一个复量
\begin{equation*}
\alpha(\omega) = \alpha'(\omega) + i\alpha''(\omega)\tag{8.21}
\end{equation*}
以描述 $D$ 与 $F$ 之间的相位差。

但 (8.20) 只不过是
\begin{equation*}
D(t) = \int_{-\infty}^{\infty} A(t-t^{\prime})\,F(t^{\prime})\,dt^{\prime},\tag{8.22}
\end{equation*}
的 Fourier 变换；其中时刻 $t$ 的位移是系统对在所有其他时刻 $t^{\prime}$ 作用的力的\emph{线性响应}（linear response）的总和。但是在力施加\emph{之后}才可能有位移：响应函数受到\emph{因果性}（causality）这一严格条件的约束：
\begin{equation*}
A(t-t^{\prime}) \equiv 0 \qquad \text{当 } t^{\prime} > t \text{ 时}.\tag{8.23}
\end{equation*}
换言之，极化率的 Fourier 积分被完全确定为
\begin{equation*}
\alpha(\omega) = \int_{-\infty}^{\infty} A(t)\,e^{i\omega t}\,dt \equiv \int_0^{\infty} A(t)\,e^{i\omega t}\,dt.\tag{8.24}
\end{equation*}
其中不包含任何来自 $t$ 负值的贡献。

为赋予这类积分以意义，通常把 $\alpha(\omega)$ 当作\emph{复}变量 $\omega$ 的函数来处理，比如说让它带上一个无穷小的正虚部 $i\epsilon$。此外，(8.24) 的形式还意味着：这个函数在复 $\omega$ 平面上实轴上方没有奇点，并且在下半平面的所有方向上随 $|\omega|\to\infty$ 而趋于零。稍作一点泛函分析——把 Cauchy 定理应用于一条沿实轴行进、再经无穷大半圆返回的围道——便得到公式
\begin{equation*}
\alpha(\Omega) = \frac{1}{i\pi}\int_{-\infty}^{\infty}\frac{\alpha(\omega)}{\omega-\Omega}\,d\omega\tag{8.25}
\end{equation*}
其中积分取其主值，而频率变量为实数。

关于“负频率”的 Fourier 变换的定义意味着 $\alpha(-\omega)$ 是 $\alpha(\omega)$ 的复共轭。由 (8.25) 我们推出极化率 (8.21) 的实部与虚部之间的 \emph{Kramers--Kronig 关系}（Kramers--Kronig relations）
\begin{equation*}
\alpha'(\Omega) = \frac{2}{\pi}\int_0^{\infty}\frac{\omega\,\alpha''(\omega)\,d\omega}{\omega^2-\Omega^2} + \mathrm{const.},\tag{8.26}
\end{equation*}
以及
\begin{equation*}
\alpha''(\Omega) = -\frac{2\Omega}{\pi}\int_0^{\infty}\frac{\alpha'(\omega)\,d\omega}{\omega^2-\Omega^2}\,.\tag{8.27}
\end{equation*}
在我们这里的情形，$\mathsf{N}^2$ 的实部和虚部分别是
\begin{equation*}
\mathsf{N}^2 = (\mathsf{n}^2 - \mathsf{k}^2) + 2\mathsf{n}\mathsf{k}i,\tag{8.28}
\end{equation*}
于是我们得到如下条件
\begin{equation*}
\mathsf{n}^2(\omega) - \mathsf{k}^2(\omega) = \frac{2}{\pi}\int_0^{\infty}\frac{\omega'\,2\mathsf{n}(\omega^{\prime})\,\mathsf{k}(\omega^{\prime})\,d\omega^{\prime}}{\omega'^{2}-\omega^{2}} + \mathrm{const.},\tag{8.29}
\end{equation*}
以及
\begin{equation*}
2\mathsf{n}(\omega)\,\mathsf{k}(\omega) = -\frac{2\omega}{\pi}\int_0^{\infty}\frac{\left\{\mathsf{n}^2(\omega^{\prime}) - \mathsf{k}^2(\omega^{\prime})\right\}d\omega^{\prime}}{\omega'^{2}-\omega^{2}}\,.\tag{8.30}
\end{equation*}
由这些关系，假如我们比如说知道了吸收系数在所有频率上作为频率的函数，那么就能分别求出 $\mathsf{n}(\omega)$ 和 $\mathsf{k}(\omega)$。它们并不是某个特定模型的偶然结果；其证明仅仅依赖于因果性条件 (8.23)，是完全普遍的，适用于任何线性响应函数，例如介电函数（dielectric function）(5.16) 或表面阻抗 (8.106)。

\section{色散与吸收}

固体的最简单模型，是一组固定不动的、彼此独立的中性原子。电磁波对这样的系统会有什么影响？倘若读者对这一理论不太熟悉，我们就在一个简单的例子中把它勾画出来：设每个原子只含一个电子，该电子处于基态轨道 $\phi_0(\mathbf{r})$ 中，可以被激发到较高的轨道 $\phi_j(\mathbf{r})$。

设波中原子附近的电场由下式给出
\begin{equation*}
\mathbf{E}(t) = \mathbf{E}_x\left(e^{i\omega t} + e^{-i\omega t}\right),\tag{8.31}
\end{equation*}
其中忽略了随距离的变化，正如 (8.3) 中那样。我们采用含时微扰论，并设电子波函数在任一特定时刻可以写成
\begin{equation*}
\psi(\mathbf{r},t) = \phi_0\exp(-i\mathcal{E}_0 t/\hbar) + \sum_j c_j(t)\,\phi_j\exp(-i\mathcal{E}_j t/\hbar).\tag{8.32}
\end{equation*}
我们把 $e\mathbf{E}\cdot\mathbf{r}$ 当作作用在电子上的微扰势。把 (8.32) 代入含时 Schrödinger 方程，用 $\phi_j^{*}$ 乘遍各项再积分（译注：原文为“用 $\phi_0$ 相乘”，按式 (8.33) 应为 $\phi_j^{*}$），容易证明：较高态的系数必须近似满足微分方程
\begin{equation*}
i\hbar\frac{dc_j}{dt} = \int \phi_j^{*}\,e\mathbf{E}(t)\cdot\mathbf{r}\,\phi_0\,d\mathbf{r}\,e^{i(\mathcal{E}_j-\mathcal{E}_0)t/\hbar}\tag{8.33}
\end{equation*}
它有解
\begin{align*}
c_j(t) &= \frac{1}{i\hbar}\int_0^t e\mathbf{E}_x\,x_{j0}\left(e^{i\omega t}+e^{-i\omega t}\right)e^{i(\mathcal{E}_j-\mathcal{E}_0)t/\hbar}\\
&= e\mathbf{E}_x\,x_{j0}\left\{\frac{1-e^{i(\hbar\omega+\mathcal{E}_j-\mathcal{E}_0)t/\hbar}}{\hbar\omega+(\mathcal{E}_j-\mathcal{E}_0)} - \frac{1-e^{i(-\hbar\omega+\mathcal{E}_j-\mathcal{E}_0)t/\hbar}}{\hbar\omega-(\mathcal{E}_j-\mathcal{E}_0)}\right\},\tag{8.34}
\end{align*}
其中
\begin{equation*}
e\,x_{j0} \equiv \int \phi_j^{*}\,e\,x\,\phi_0\,d\mathbf{r},\tag{8.35}
\end{equation*}
也就是说，它是电子的电偶极矩沿场的电矢量方向、在态 $\phi_0$ 与 $\phi_j$ 之间的矩阵元。

现在，量子理论书中的通常做法是研究 $|c_j(t)|^2$，并证明：除了满足条件
\begin{equation*}
\hbar\omega = \hbar\omega_j \equiv (\mathcal{E}_j - \mathcal{E}_0).\tag{8.36}
\end{equation*}
的那些 $\omega$ 值之外，它都剧烈振荡。于是人们证明，将会发生一次从态 $\phi_0$ 到态 $\phi_j$ 的量子跃迁，其概率由矩阵元 (8.35) 的平方决定。

但现在我们改走另一条路：计算原子在这个含时态中的偶极矩之值。由 (8.32) 和 (8.34) 我们得到
\begin{align*}
\langle ex(t)\rangle &= \int \psi^{*}(\mathbf{r},t)\,ex\,\psi(\mathbf{r},t)\,d\mathbf{r}\\
&= \sum_j \left\{ex_{0j}\,c_j(t)\,e^{-i\omega_j t} + ex_{j0}\,c_j^{*}(t)\,e^{i\omega_j t}\right\}\\
&= \mathbf{E}_x\sum_j e^2\,|x_{0j}|^2\,\frac{1}{\hbar}\left\{\frac{1}{\omega_j-\omega}+\frac{1}{\omega_j+\omega}\right\}\left(e^{i\omega t}+e^{-i\omega t}\right)\tag{8.37}
\end{align*}
以及一些振荡迅速、与应用场不同相的项。

这样，原子的偶极矩正比于场。我们说\emph{原子极化率}（atomic polarizability）由下式给出
\begin{equation*}
\alpha(\omega) = \sum_j \frac{e^2|x_{0j}|^2}{\hbar}\,\frac{2\omega_j}{\omega_j^2-\omega^2}\,.\tag{8.38}
\end{equation*}
为讨论这类表达式的数量级，我们可以利用 Thomas--Reiche--Kuhn 求和规则——它是 (5.48) 的一个较简单的特例。这个规则是
\begin{equation*}
\frac{2m}{\hbar^2}\sum_j \hbar\omega_j\,|x_{0j}|^2 = 1\,;\tag{8.39}
\end{equation*}
为方便起见，定义第 $j$ 个跃迁的\emph{振子强度}（oscillator strength）
\begin{equation*}
f_j = \frac{2m}{\hbar^2}\,\hbar\omega_j\,|x_{0j}|^2,\tag{8.40}
\end{equation*}
并把 (8.38) 写成如下形式
\begin{equation*}
\alpha(\omega) = \frac{e^2}{m}\sum_j \frac{f_j}{\omega_j^2-\omega^2},\tag{8.41}
\end{equation*}
其中各个数 $f_j$ 在所有跃迁上相加为一。

现在，若把 $N$ 个这样的原子聚拢到每单位体积中，我们就得到一种介质，其介电常数有一个实部
\begin{align*}
\epsilon(\omega) &= 1 + 4\pi N\alpha(\omega)\\
&= 1 + \frac{4\pi N e^2}{m}\sum_j \frac{f_j}{\omega_j^2-\omega^2}\,.\tag{8.42}
\end{align*}
严格说来，我们应当把这个\emph{横}（transverse）介电常数，与第 5 章中讨论过的静电场的\emph{纵}（longitudinally）偏振波的介电常数区分开来。不过在长波极限下——即对通常的“光学”现象——这两个量是相同的，因为屏蔽或原子极化效应相对地说是局域的，并且主要依赖于每个离子或原子紧邻区域中电场的强度。
这是色散公式的原型。有趣的是，它恰好可以归结为 \S 5.7 中已经讨论过的情形。如果电子是“自由的”——也就是说，所有 $\omega_j$ 均为零，而各个 $f_j$ 之和为一——我们便可以像 (5.53) 那样写出
\begin{equation*}
\omega_p^2 = \frac{4\pi N e^2}{m}\tag{8.43}
\end{equation*}
来为我们这个每单位体积 $N$ 个电子的系统定义等离体频率（plasma frequency）；于是 (8.42) 就复现了 (5.67)，即导致紫外透明性现象的公式。然而要注意，横向光子实际上几乎不能直接激发纵向等离激元（plasmon）。要观察\emph{等离激元谱线}（plasmon lines），需要用高能电子去激发电子气，或者透过薄膜去观察，在后一种情形中可以观察到\emph{表面等离激元模}（surface plasmon modes）。

但在一般情况下，$\omega_j$ 的最低值应对应于光学或红外频率，$\epsilon$ 随 $\omega$ 的变化要更复杂一些。在低频下，我们可能得到静态介电常数
\begin{equation*}
\epsilon(0) = 1 + \sum_j f_j\,\frac{\omega_p^2}{\omega_j^2},\tag{8.44}
\end{equation*}
它明显大于 1。但 $\epsilon(0)$ 不会非常大，因为对大多数固体来说，$\omega_p$ 与自由原子的普通光学频率大致相当。

接着，随着 $\omega$ 增大，$\epsilon(\omega)$ 增大，直到在 $\omega = \omega_j$ 处出现一个奇点。此后，$\epsilon(\omega)$ 在某段频率范围内变为负值，直到靠近下一个共振频率时才重新变正。不过最终，$\epsilon(\omega)$ 将趋于公式 (5.67)：
\begin{equation*}
\epsilon(\omega) \to 1 - \frac{\omega_p^2}{\omega^2},\tag{8.45}
\end{equation*}
条件是 $\omega$ 大于所有 $\omega_j$。

在每一个 $\omega_j$ 附近，我们都有一段\emph{反常色散}（anomalous dispersion）区。特别地，若 $\epsilon$ 为负，则由 (8.7)，折射率 $\mathsf{N}$ 变为纯虚数；于是
\begin{equation*}
\mathsf{n} = 0,\quad \mathsf{k} = |\epsilon|^{\frac{1}{2}},\tag{8.46}
\end{equation*}
从而由 (8.19)，光将被固体表面\emph{完全反射}（total reflection）。实际上，晶体看起来会一直保持透明——尽管折射率非常高——直到 $\omega = \omega_j$ 处突然变得不透明且完全反射，然后在更高的频率上重新变为透明。

但这与我们的色散关系 (8.29) 和 (8.30) 不相容。还必须存在某种吸收，其量度是乘积 $2\mathsf{n}\mathsf{k}\omega$。的确，这正是我们在含时微扰处理中暂且搁置的那一方面——与原子跃迁 (8.36) 相联系的吸收。它作为频率 $\omega_j$ 处的一条锐线出现——想必是一个 delta 函数。为了满足色散关系 (8.29) 并在 (8.42) 中给出正确的色散项，我们必须有
\begin{equation*}
2\omega\,\mathsf{n}(\omega)\,\mathsf{k}(\omega) = \frac{\pi}{2}\,\frac{4\pi N e^2}{m}\sum_j f_j\,\delta(\omega-\omega_j);\tag{8.47}
\end{equation*}
%FIG{137}: {光的色散。}
这必定是复介电常数的虚部，后者取如下形式
\begin{equation*}
\epsilon(\omega) = 1 + \frac{4\pi N e^2}{m}\sum_j f_j\left\{\frac{1}{\omega_j^2-\omega^2} + \frac{i\pi}{2\omega}\,\delta(\omega^2-\omega_j^2)\right\}.\tag{8.48}
\end{equation*}
$\omega_j$ 附近的大折射率在这个频率上变成一条尖锐的吸收线。

实际上，我们永远不会看到一条无限锐的谱线。总会有某种展宽，来自杂质等因素，或者仅仅来自能级的自然辐射弛豫。正如标准量子理论所表明的，这类效应可以在分析中以唯象方式纳入：比如说在 (8.34) 中插入一个衰减因子 $\exp(-\frac{1}{2}\Gamma t)$，它对应于（振幅\emph{平方}的）量级为 $1/\Gamma$ 秒的衰减时间。

在代数推导中很容易追踪这样一个因子的作用：它相当于
%FIG{138}: {折射率的实部和虚部。}
%FIG{139}: {色散曲线的展宽。}
在 (8.37) 或 (8.38) 的 $\omega_j$ 上加一个 $\frac{1}{2}i\Gamma$。如果相对于 $\omega_j^2$ 略去 $\Gamma^2$，那么对于极化率的实部——从而也就是 $\epsilon(\omega)$ 的实部——我们会得到如下类型的项
\begin{equation*}
f_j\,\frac{\omega_j^2-\omega^2}{(\omega_j^2-\omega^2)^2+\omega^2\Gamma^2}\tag{8.49}
\end{equation*}
用以取代 $f_j/(\omega_j^2-\omega^2)$。这起到了消除 $\mathsf{n}$ 中奇点的作用，
并把色散函数在频率上铺展开来，形成宽度约为 $2\Gamma$ 的一段范围。

众所周知，弛豫对吸收线本身的影响，是把 delta 函数铺展成一个有限大小的函数，在 $\omega = \omega_j$ 附近其形式为
\begin{equation*}
\frac{\Gamma/2\pi}{(\omega_j-\omega)^2+\frac{1}{4}\Gamma^2} \approx \frac{2\Gamma\omega^2/\pi}{(\omega_j^2-\omega^2)^2+\Gamma^2\omega^2}\tag{8.50}
\end{equation*}
%FIG{140}: {Lorentz 修正：用均匀极化介质中一个球形空腔的场，取代相邻原子上的偶极子在原子处产生的局域场。}
实际上，(8.49) 与 (8.50) 可以合并成一个类似 (8.48) 的公式
\begin{align*}
\epsilon(\omega) &= 1 + \frac{4\pi N e^2}{m}\sum_j f_j\left\{\frac{\omega_j^2-\omega^2}{(\omega_j^2-\omega^2)^2+\omega^2\Gamma_j^2} + \frac{i\Gamma_j\omega}{(\omega_j^2-\omega^2)^2+\omega^2\Gamma_j^2}\right\}\\
&= 1 + \frac{4\pi N e^2}{m}\sum_j \frac{f_j}{(\omega_j^2-\omega^2)-i\omega\Gamma_j},\tag{8.51}
\end{align*}
它差不多就是最一般类型的色散公式。虽然它是针对独立原子模型算出的，却为任何吸收谱由一系列分立谱线组成的系统提供了一个唯象表达式。

在这个分析中，还应作一个进一步的修正。关系 (8.42) 意味着极化每个原子的\emph{局域}场（local field）与施加于晶体的\emph{宏观}场（macroscopic field）相同。严格说来这并不成立；原子并非处于一种绝对连续的
%JOIN%
